\clearpage
\begin{trapbox}
\textbf{\textcolor{CTrap}{易错一、二：丢绝对值，或误把 $F_1$ 永远当作远焦点}}

原式是 $d_1=|a+ex_0|,d_2=|a-ex_0|$。先确认点在曲线上，再按所在支判断括号的正负。
\begin{center}
\begin{tikzpicture}[x=1cm,y=1cm,line cap=round,line join=round,font=\small]
\draw[black!25,rounded corners=3pt] (0,0) rectangle (14.4,9.5);
\node[font=\small\bfseries] at (7.2,9.0) {先看允许范围，再看括号正负：绝对值不能直接丢};
\node at (7.2,8.15) {$e>1\ \Longrightarrow\ 0<a/e<a$；两个零点都落在无曲线点的中间区域。};
\fill[blue!8] (1,5.85) rectangle (4.8,6.55);
\fill[black!7] (4.8,5.85) rectangle (9.6,6.55);
\fill[orange!10] (9.6,5.85) rectangle (13.4,6.55);
\draw[->,black!70] (1,6.2)--(13.5,6.2) node[right] {$x_0$};
\draw[line width=1.5pt,blue!65] (1.1,6.2)--(4.8,6.2);
\draw[line width=1.5pt,orange!85!black] (9.6,6.2)--(13.2,6.2);
\fill (4.8,6.2) circle(2pt); \fill (9.6,6.2) circle(2pt);
\node[above] at (4.8,6.55) {$-a$}; \node[above] at (9.6,6.55) {$a$};
\draw[purple!70] (5.6,6.05)--(5.6,6.35) (8.8,6.05)--(8.8,6.35);
\node[below] at (5.6,5.8) {$-\dfrac ae$}; \node[below] at (8.8,5.8) {$\dfrac ae$};
\node at (7.2,4.55) {$-a<x_0<a$：没有双曲线上的点};
\node at (2.8,7.45) {左支 $x_0\le-a$}; \node at (11.6,7.45) {右支 $x_0\ge a$};
\node[align=center] at (2.8,3.35) {$a+ex_0<0$\\[.6em]$a-ex_0>0$};
\node[align=center] at (11.6,3.35) {$a+ex_0>0$\\[.6em]$a-ex_0<0$};
\node[align=center] at (2.8,1.55) {\textcolor{blue!70!black}{$d_1=-ex_0-a$}\\[.6em]\textcolor{orange!85!black}{$d_2=-ex_0+a$}};
\node[align=center] at (11.6,1.55) {\textcolor{blue!70!black}{$d_1=ex_0+a$}\\[.6em]\textcolor{orange!85!black}{$d_2=ex_0-a$}};
\end{tikzpicture}
\end{center}

\textbf{反例：} 主图右支点 $x_0=3,a=2,e=3/2$，
若错用 $d_2=a-ex_0$，会得到 $-5/2$。距离不可能为负，应取绝对值得 $5/2$。
左支则由 $F_1$ 变成同侧近焦点（见左右支对照图），焦点下标不能按远近重新命名。
\end{trapbox}
\clearpage
\begin{trapbox}
\textbf{\textcolor{CTrap}{易错三：焦点换到纵轴，仍使用横坐标}}
\begin{center}
\begin{tikzpicture}[x=1cm,y=1cm,line cap=round,line join=round,font=\small]
\draw[black!25,rounded corners=3pt] (0,0) rectangle (14.4,10);
\node[font=\small\bfseries] at (7.2,9.5) {焦点轴变成纵轴：坐标与分支一起换成 $y_0$};
\begin{scope}[shift={(3.6,5.3)},x=0.6cm,y=0.6cm]
\draw[black!30,->] (-4.4,0)--(4.5,0) node[right] {$x$};
\draw[black!30,->] (0,-3.8)--(0,4.1) node[above] {$y$};
\draw[black!65,thick,domain=-3.5:3.5,samples=100,variable=\z] plot ({-1*2*sqrt(1+\z*\z/5)},{\z});
\draw[black!65,thick,domain=-3.5:3.5,samples=100,variable=\z] plot ({1*2*sqrt(1+\z*\z/5)},{\z});
\coordinate (P) at (3,2.5);
\coordinate (Fone) at (-3,0);
\coordinate (Ftwo) at (3,0);
\draw[blue!70!black,line width=1.1pt] (Fone)--(P);
\draw[orange!85!black,line width=1.1pt] (Ftwo)--(P);
\foreach \q in {P,Fone,Ftwo} \fill (\q) circle (1.6pt);
\node[above left] at (P) {$P$};
\node[below left] at (Fone) {$F_1$};
\node[below right] at (Ftwo) {$F_2$};
\end{scope}
\node at (3.6,8.9) {左、右焦点；图示右支};
\node[blue!70!black] at (3.6,1.9) {$d_1=|a+ex_0|$};
\node[orange!85!black] at (3.6,1.2) {$d_2=|a-ex_0|$};
\begin{scope}[shift={(10.8,5.3)},x=0.6cm,y=0.6cm]
\draw[black!30,->] (-4.4,0)--(4.5,0) node[right] {$x$};
\draw[black!30,->] (0,-3.8)--(0,4.1) node[above] {$y$};
\draw[black!65,thick,domain=-3.5:3.5,samples=100,variable=\z] plot ({\z},{-1*2*sqrt(1+\z*\z/5)});
\draw[black!65,thick,domain=-3.5:3.5,samples=100,variable=\z] plot ({\z},{1*2*sqrt(1+\z*\z/5)});
\coordinate (P) at (-2.5,3);
\coordinate (Fone) at (0,-3);
\coordinate (Ftwo) at (0,3);
\draw[blue!70!black,line width=1.1pt] (Fone)--(P);
\draw[orange!85!black,line width=1.1pt] (Ftwo)--(P);
\foreach \q in {P,Fone,Ftwo} \fill (\q) circle (1.6pt);
\node[above right] at (P) {$P$};
\node[below left] at (Fone) {$F_1$};
\node[below right] at (Ftwo) {$F_2$};
\end{scope}
\node at (10.8,8.9) {下、上焦点；图示上支};
\node[blue!70!black] at (10.8,1.9) {$d_1=|a+ey_0|$};
\node[orange!85!black] at (10.8,1.2) {$d_2=|a-ey_0|$};
\node at (7.2,.5) {两图等比例、同尺度，右图由左图旋转 $90^\circ$ 得到，两条焦半径长度不变。};
\end{tikzpicture}
\end{center}

右支图的 $P=(3,5/2)$ 旋转后成为上支点 $P'=(-5/2,3)$。
此时 $F_1$ 是下焦点、$F_2$ 是上焦点，公式中的横坐标改成纵坐标。
本图两种状态的 $d_1,d_2$ 都为 $13/2,5/2$。
\end{trapbox}
\clearpage
\begin{trapbox}
\textbf{\textcolor{CTrap}{易错四：只看横坐标符号，或只检查 $|x_0|\ge a$}}
例如 $a=4$ 时 $x_0=2>0$，仍不可能是双曲线点；
在 $x^2/16-y^2/9=1$ 中，$(5,0)$ 虽满足 $x_0\ge4$，却不在曲线上。
\textbf{纠正：} 范围只是必要条件，必须核对原方程。

\medskip
\textbf{\textcolor{CTrap}{易错五：把 $c^2$ 写成 $a^2-b^2$}}
双曲线始终是 $c^2=a^2+b^2$，因而 $e>1$；先认清曲线，再算参数。

\medskip
\textbf{\textcolor{CTrap}{易错六：平方开方直接丢掉绝对值}}
$\sqrt{(ex_0-a)^2}=|ex_0-a|$。先保留绝对值，或用前面的数轴图判断符号。

\medskip
\textbf{\textcolor{CTrap}{易错七：把椭圆固定和的设法搬来}}
椭圆和为 $2a$，可设 $a+t,a-t$；双曲线差的绝对值为 $2a$，应设 $t+a,t-a$。
用方法二的实长线段图检查：两端分别距平均长度 $a$，总差才是 $2a$。

\medskip
\textbf{\textcolor{CTrap}{易错八：把平均长度 $t$ 当成原图中的几何中点}}
$t=(d_1+d_2)/2$ 是\textbf{两个数的平均值}。原图中两线段不共线，不能据此断言某个几何中点到焦点距离为 $t$。

\medskip
\textbf{考场回验：} 长减短为 $2a$，长加短为 $2e|x_0|$；顶点也应满足这两式。
\end{trapbox}
